\documentclass{tufte-handout}

\usepackage{style}

\begin{document}

\tma{04} % TMA number

\section{Part A}

%Question 1
\begin{question}

    Let \( R \) be the ring,
    \[ \Z[\sqrt{-14}] \]

    \qpart

    Let \( A \) be the subring,
    \[ A = \{ a + b\sqrt{-14} : a,b \in \Z, a - b \equiv 0\pmod{5} \} \]

    As \( a = b\pmod{5} \) there exists some \( k \in \Z \) such that \( a = b + 5k \).

    Hence,
    \begin{align*}
    a + b\sqrt{-14} &= b + 5k + b\sqrt{-14}\\[8pt]
    &= b(1 + \sqrt{-14}) + 5k\\[8pt]
    \stext{So}
    A &= \langle 1 + \sqrt{-14}, 5 \rangle
    \end{align*}

    By \textup{Definition 2.1 HB chapter 19 p94}, and \( R = \Z[\sqrt{-14}] = \langle 1, \sqrt{-14} \rangle_{\Z} \),
    so we only need to check multipliation by \( \sqrt{-14} \), as \( 1 \) is trival.

    \begin{align*}
    \sqrt{-14}(1 + \sqrt{-14}) &= \sqrt{-14} + (\sqrt{-14})^2\\[8pt]
    \stext{and since}
    -14 \equiv 1 \pmod{5} \\[8pt]
    \stext{and}
    -14 = 1 -15\\[8pt]
    \stext{So}
    \sqrt{-14} - 14 &= (\sqrt{-14} + 1) - 15\\[8pt]
    &= (\sqrt{-14} + 1) + 5(-3)\\[8pt]
    \end{align*}

    AS \( \sqrt{-14} \in A \) and \( 15 = 3(5) \in A \), it follows that \( A \) is closed under multiplication by any 
    element of \( R \), and hence \( A \) is an ideal of \( R \).

    \qpart



\end{question}

%Question 2
\begin{question}

\[ f(x) = 3x^{4} + 156x^{3} + 26x^{2} +130x +26 \in Z[x] \]

\qpart

Using \textup{Thereom 4.15 HB Chapter 12 p52},

using the prime \( p = 13 \)
\begin{itemize}
\item \( 13 \mid 156 \)\\
\item \( 13 \mid 26 \)\\
\item \( 13 \mid 130 \)\\
\item \( 13 \mid 26 \)\\
\item \( 13 \nmid 3 \)\\
\item \( 13^2 \nmid 26 \)\\
\end{itemize}

Hence by Eisenstein's criterion, \( f(x) \) is irreducible over \( \Q[x] \).

\vspace{2cm}

\qpart

Let,
\[ \alpha \in \C \quad \beta \in \C \]
be two distinct roots of \( f(x) \)

Since
\[ f(x) \in \Z[x] \subset \Q[x] \]
is irreducible over \( \Q \), and \( \alpha \in \C \) satisfies \( f(\alpha) = 0 \)
it follows that the minimal polynomial of \( \alpha \) over \( \Q \) is \( f(x) \).

Using \textup{Theorem 2.21 HB Chapter 18 p105},
\[ \Q(\alpha) \cong \Q[x]/(f(x)) \]

The same applies to \( \beta \), and hence
\[ \Q(\beta) \cong \Q[x]/(f(x)) \]

Thus,
\[ \Q(\alpha) \cong \Q(\beta) \]

\end{question}

%Question 3
\begin{question}

    Let
    \[ \alpha = \sqrt{13} + \sqrt{n} \]

    \qpart

    \begin{align*}
    \alpha^{3} - (n + 39)\alpha &= (\sqrt{13} + \sqrt{n})^{3} - (n + 39)(\sqrt{13} + \sqrt{n})\\[8pt]
    &= (\sqrt{13} + \sqrt{n})(\sqrt{13} + \sqrt{n})^2 - (n + 39)(\sqrt{13} + \sqrt{n})\\[8pt]
    &= (\sqrt{13} + \sqrt{n})(13 + n + 2\sqrt{13}\sqrt{n}) - (n + 39)(\sqrt{13} + \sqrt{n})\\[8pt]
    &= (\sqrt{13} + \sqrt{n})((13 + n + 2\sqrt{13n}) - (n + 39))\\[8pt]
    &= (\sqrt{13} + \sqrt{n})(-26 + 2\sqrt{13n})\\[8pt]
    &= -26\sqrt{13} + 2\cdot\sqrt{13}\sqrt{13n} - 26\sqrt{n} + 2\sqrt{13n^2}\\[8pt]
    &= -26\sqrt{13} + 2\cdot13\sqrt{n} - 26\sqrt{n} + 2\sqrt{13}n\\[8pt]
    &= -26\sqrt{13} + 26\sqrt{n} - 26\sqrt{n} + 2\sqrt{13}n\\[8pt]
    &= -26\sqrt{13} + 2\sqrt{13}n\\[8pt]
    &= \sqrt{13}(2n - 26)\\[8pt]
    \end{align*}

    Hence,
    \[ \alpha^{3} - (n + 39)\alpha = \sqrt{13}(2n - 26) \]

    and we can write this as,
    \[ \sqrt{13} = \frac{\alpha^{3} - (n + 39)\alpha}{(2n - 26)} \]

    So
    \[ \sqrt{13} \in \Q_{\alpha} \]

    and since \( \sqrt{n} = \alpha - \sqrt{13} \), it follows that
    \[ \sqrt{n} \in \Q_{\alpha} \]

    Hence,
    \[ \Q(\sqrt{13}, \sqrt{n}) \subseteq \Q_{\alpha} \]

    Hence,
    \[ \Q(\sqrt{13}, \sqrt{n}) \subseteq \Q(\sqrt{13} + \sqrt{n}) \]

\end{question}


%Question 4
\begin{question}

    Let
    \[ /bm{F}  =  \]

\end{question}


\end{document}